\subsection{SPLIT SEMI-SIMPLE LIE ALGEBRAS}

DEFINITION 1. Let \(\mathfrak{g}\) be a semi-simple Lie algebra. A Cartan subalgebra \(\mathfrak{h}\) of \(\mathfrak{g}\) is called splitting if, for all \(x \in \mathfrak{h}\) , \(\operatorname{ad}_{\mathfrak{g}} x\) is triangularizable. A semi-simple Lie algebra is called splittable if it has a splitting Cartan subalgebra. A split semi-simple Lie algebra is a pair \((\mathfrak{g}, \mathfrak{h})\) where \(\mathfrak{g}\) is a semi-simple Lie algebra and \(\mathfrak{h}\) is a splitting Cartan subalgebra of \(\mathfrak{g}\) .

Remarks. 1) Let \(\mathfrak{g}\) be a semi-simple Lie algebra, \(\mathfrak{h}\) a Cartan subalgebra of \(\mathfrak{g}\) . For all \(x \in \mathfrak{h}\) , \(\operatorname{ad}_{\mathfrak{g}} x\) is semi-simple (Chap. VII, §2, no. 4, Th. 2). Thus, to say that \(\mathfrak{h}\) is splitting means that \(\operatorname{ad}_{\mathfrak{g}} x\) is diagonalizable for all \(x \in \mathfrak{h}\) .

\begin{enumerate}
\def\labelenumi{\arabic{enumi})}
\setcounter{enumi}{1}
\item
  If \(k\) is algebraically closed, every semi-simple Lie algebra \(\mathfrak{g}\) is splittable, and every Cartan subalgebra of \(\mathfrak{g}\) is splitting. When \(k\) is not algebraically closed, there exist non-splittable semi-simple Lie algebras (Exerc. 2 a)); moreover, if \(\mathfrak{g}\) is splittable, there may exist Cartan subalgebras of \(\mathfrak{g}\) that are not splitting (Exerc. 2 b)).
\item
  Let \(\mathfrak{g}\) be a semi-simple Lie algebra, \(\mathfrak{h}\) a Cartan subalgebra of \(\mathfrak{g}\) , and \(\rho\) a finite dimensional injective representation of \(\mathfrak{g}\) such that \(\rho(\mathfrak{h})\) is diagonalizable. Then \(\operatorname{ad}_{\mathfrak{g}}x\) is diagonalizable for all \(x \in \mathfrak{h}\) (Chap. VII, §2, no. 1, Example 2), so \(\mathfrak{h}\) is splitting.
\item
  We shall see (§3, no. 3, Cor. of Prop. 10) that if \(\mathfrak{h}\) , \(\mathfrak{h}'\) are splitting Cartan subalgebras of \(\mathfrak{g}\) , there exists an elementary automorphism of \(\mathfrak{g}\) transforming \(\mathfrak{h}\) into \(\mathfrak{h}'\) .
\item
  Let \(\mathfrak{g}\) be a reductive Lie algebra. Then \(\mathfrak{g} = \mathfrak{c} \times \mathfrak{s}\) where \(\mathfrak{c}\) is the centre of \(\mathfrak{g}\) and \(\mathfrak{s} = \mathcal{D}\mathfrak{g}\) is semi-simple. The Cartan subalgebras of \(\mathfrak{g}\) are the subalgebras of the form \(\mathfrak{h} = \mathfrak{c} \times \mathfrak{h}'\) where \(\mathfrak{h}'\) is a Cartan subalgebra of \(\mathfrak{s}\) (Chap. VII, §2, no. 1, Prop. 2). Then \(\mathfrak{h}\) is called splitting if \(\mathfrak{h}'\) is splitting relative to \(\mathfrak{s}\) . This leads in an obvious way to the definition of splittable or split reductive algebras.
\end{enumerate}
% semantic-fragments: legacy-input-seam
\relax{}
